% Auriga theme
% https://github.com/anishathalye/auriga

\documentclass[14pt,aspectratio=169]{beamer}
\usepackage{pgfpages}
\usepackage{fancyvrb}
\usepackage{booktabs,dcolumn}
\usepackage{tikz}
\usetikzlibrary{arrows.meta, positioning, quotes}
\usepackage{pgfplots}

% \ifnotes
% \setbeamertemplate{note page}[plain]
% \setbeameroption{show notes on second screen=right}
% \fi

\usetheme{auriga}
\usecolortheme{auriga}

% define some colors for a consistent theme across slides
\definecolor{red}{RGB}{181, 23, 0}
\definecolor{blue}{RGB}{0, 118, 186}
\definecolor{gray}{RGB}{146, 146, 146}

\title{Historia Económica Argentina}

\author{\underline{Sebastián Freille} \inst{1}}

\institute[shortinst]{\inst{1} IEF-UNC \\ UCC}
\date{}

\AtBeginSection[]{
    \begin{frame}
    \vfill
    \centering
    \begin{beamercolorbox}[sep=8pt,center,shadow=true,rounded=true]{title}
        \usebeamerfont{title}\insertsectionhead\par%
    \end{beamercolorbox}
    \vfill
    \end{frame}
}


\AtBeginSubsection[]{
    \begin{frame}
    \vfill
    \centering
    \begin{beamercolorbox}[sep=6pt,center,shadow=true,rounded=true]{title}
        \usebeamerfont{title}\insertsubsectionhead\par%
    \end{beamercolorbox}
    \vfill
    \end{frame}
}


\AtBeginSubsubsection[]{
    \begin{frame}
    \vfill
    \centering
    \begin{beamercolorbox}[sep=4pt,center,shadow=true,rounded=true]{title}
        \usebeamerfont{title}\insertsubsubsectionhead\par%
    \end{beamercolorbox}
    \vfill
    \end{frame}
}



\begin{document}
\maketitle

{
  % rather than use the frame options [noframenumbering,plain], we make the
  % color match, so that the indicated page numbers match PDF page numbers
  \setbeamercolor{page number in head/foot}{fg=background canvas.bg}
  \begin{frame}
    \titlepage
  \end{frame}
}

\section{Aspectos generales y administrativos}

 \begin{frame}\frametitle{Objetivos}
   \begin{block}{Objetivo general}
     Analizar la evolución de la economía argentina durante el período
    que comienza con la consolidación del Estado Nacional en 1880
    hasta la década de 1990. Se buscará ofrecer a los alumnos la
    información necesaria sobre el funcionamiento de la economía
    argentina, así como también brindar una explicación de sus cambios
    mediante el uso de la teoría económica como base. El énfasis
    estará puesto en el examen de las políticas económicas aplicadas y
    los resultados alcanzados
    \end{block}
  \end{frame}


 \begin{frame}\frametitle{Equipo de la cátedra}
   \begin{itemize}
   \item Profesor a cargo: Dr. Sebastián Freille (Prof. Adjunto por
     concurso) --en reemplazo temporario de Prof. Titular Mónica Gómez
     (licencia por año sabático)
   \item Profesor ayudante: Lic. Lucas Tossolini
   \item Adscriptos: Lic. Martín Calvo / Lic. Fernando Kühn
          \end{itemize}
  \end{frame}
  

  \begin{frame}\frametitle{Metodología de evaluación}
     \begin{itemize}
     \item Habrá 2 (dos) evaluaciones parciales y 1 (un) parcial
       recuperatorio. Las fechas tentativas y propuestas a Dirección
       de Enseñanza: 
       \begin{itemize} \medskip
       \item Jueves 01/OCT de 17:00 a 19:00
       \item Martes 03/NOV de 17:00 a 19:00
         \item Martes 11/NOV de 17:00 a 19:00
         \end{itemize}
       \item La regularidad se obtiene aprobando 2 (dos) parciales con
         nota de 4 (50 puntos sobre 100) o superior
         \item Debido a las circunstancias excepcionales de evaluación
           no presencial, no habrá régimen de promoción indirecta
           durante este semestre
       \end{itemize}
     \end{frame}


     \begin{frame}\frametitle{Metodología de evaluación (cont.)}
     \begin{itemize}
     \item Alumnos regulares $\longrightarrow$ el examen será
       escrito y versara sobre la totalidad del programa. Contemplará
       preguntas de multiple choice y ensayos
       \item Alumnos libres $\longrightarrow$ el examen sera escrito
         especial y versará sobre la totalidad del programa. Tendrá
         una estructura similar al examen regular con algo más de
         preguntas. 
       \end{itemize}
  \end{frame}


  
  
  
\begin{frame}\frametitle{Bibliografía}
     \begin{itemize}
     \item Se compone de libros y articulos seleccionados. Los libros
       más consultados son: Gerchunoff y Llach (2007), Díaz Alejandro
       (1970), Arnaudo (1986) y Gómez (2018)
       \item El material de lectura es abundante y complejo en su
         cobertura y lenguaje. Se recomienda que las clases y
         las diapositivas sean consideradas como un esqueleto básico y
         mínimo a partir del cual encarar la lectura de los textos.
         \item Toda la bibliografía estará disponible en el aula
           virtual de la materia:
           \url{https://aulavirtual.eco.unc.edu.ar/enrol/index.php?id=597} 
       \end{itemize}
  \end{frame}


  \begin{frame}\frametitle{Bibliografía}
     \begin{itemize}
     \item Gerchunoff, P. y Llach, L. (2007). \textit{El ciclo de la
         ilusión y el desencanto. Un siglo de políticas económicas
         argentinas}, Buenos Aires, Ariel, Edición 2007
       \item Díaz Alejandro, Carlos (1970). \textit{Ensayos sobre la
           historia económica argentina}, Buenos Aires, Amorrortu.
                \item Arnaudo, Aldo (1986). \textit{Cincuenta años de
               política financiera}, Buenos Aires, El Ateneo.
             \item Gómez, Mónica (2018). \textit{Los avatares de la
                 primera Caja de Conversión Argentina y sus
                 transformación final en Banco Central, 1890-1935},
               Buenos Aires, Teseo. 
       \end{itemize}
  \end{frame}

  

\begin{frame}\frametitle{Canales de comunicación}
     \begin{itemize}
     \item Proponemos tres canales principales de comunicación
       \begin{itemize}\itemsep 10pt \medskip
       \item Aula virtual $\longrightarrow$ uso como repositorio de
         bibliografía en formato digital y contenidos de la
         materia. Algún tipo de comunicaciones entre docentes y
         alumnos
         \item Correo de catedra $\longrightarrow$
           \url{hea.fce.unc@gmail.com} a través del cual se enviarán
           comunicaciones regulares a alumnos
           \item Clases interactivas $\longrightarrow$ plataforma Zoom
         \end{itemize}
       \end{itemize}
  \end{frame}


  
   \begin{frame}\frametitle{Programa abreviado}
     \begin{enumerate}
     \item Capítulo 1. El período de progreso (1880-1914)
     \item Capítulo 2. Los años entre la dos guerras (1914-1945)
       \item Capítulo 3. La economía política del peronismo
         (1946-1955)
         \item Capítulo 4. Hacia la liberalización de los mercados en
           los años que siguieron al peronismo (1955-1963)
         \item Capítulo 5. La década de expansión (1963-1973)
           \item Capítulo 6. El período de estancamiento y declinación
             (1973-1989)
             \item Capítulo 7. La década de los 90. 
       \end{enumerate}
  \end{frame}


  \begin{frame}\frametitle{Programa abreviado: Capítulo 1}
     \begin{enumerate}
     \item Capítulo 1. El período de progreso (1880-1914)
       \begin{enumerate}
       \item De Roca a Saenz Peña
         \end{enumerate}
       \item Temas
         \begin{enumerate}
           \item Crecimiento económico a través de la teoría del bien
             primario exportable (BPE) [``staple theory'']
           \item Funcionamiento monetario y crisis de 1890-91
             \item Nuevo modelo monetario: el sistema de Caja de Conversión
             \end{enumerate}
             \end{enumerate}
  \end{frame}


\begin{frame}\frametitle{Programa abreviado: Capítulo 2}
     \begin{enumerate}
     \item Capítulo 2. Los años entre las dos guerras (1914-1945)
       \begin{enumerate}
       \item Los últimos años del Partido Autonomista Nacional
       \item Los gobiernos radicales
       \item La Década infame
         \item La Revolución del 43
         \end{enumerate}
       \item Temas
         \begin{enumerate}
           \item Efectos de la primera GM en la economía y las
             respuestas de politica economica
           \item Cambios en la estructura en la década del 20
             (industria)
             \item Impactos de la Gran Depresión en la economía y
               respuestas de politica economica
             \end{enumerate}
             \end{enumerate}
  \end{frame}


\begin{frame}\frametitle{Programa abreviado: Capítulo 3}
     \begin{enumerate}
     \item Capítulo 3. La economía política del peronismo (1946-1955)
       \begin{enumerate}
       \item El primer gobierno de Perón (1946-1952)
       \item El segundo gobierno de Perón (1952-1955)
       \end{enumerate}
     \item Temas
       \begin{enumerate}
         \item La economía política del peronismo y la política
           económica (1946-1949)
         \item El origen de los problemas: ``stop-go'' e inflación
           \item Cambios en la politica económica (1952-1955)
           \end{enumerate}
           \end{enumerate}
  \end{frame}


  
\begin{frame}\frametitle{Programa abreviado: Capítulo 4}
     \begin{enumerate}
     \item Capítulo 4. Hacia la liberalización de los mercados en los
       años que siguieron al peronismo (1955-1963)
       \begin{enumerate}
       \item La Revolucion Libertadora
       \item El régimen de Frondizi
         \item El gobierno de Guido
       \end{enumerate}
     \item Temas
       \begin{enumerate}
       \item El plan de Prebisch y la política económica
       \item La estrategia desarrollista de Frondizi
         \item El programa económico de 1959-63 y la crisis del 1962-1963
         \end{enumerate}
         \end{enumerate}
  \end{frame}


  \begin{frame}\frametitle{Programa abreviado: Capítulo 5}
     \begin{enumerate}
     \item Capítulo 5. La década de expansión (1963-1973)
       \begin{enumerate}
       \item El gobierno de Illia
       \item La Revolución Argentina: El gobierno de Onganía
      \end{enumerate}
     \item Temas
       \begin{enumerate}
       \item El crecimiento económico en la década de los 60
       \item La política económica durante el gobierno de Illia
         \item El plan de Krieger Vasena
         \end{enumerate}
         \end{enumerate}
  \end{frame}
  

 \begin{frame}\frametitle{Programa abreviado: Capítulo 6}
     \begin{enumerate}
     \item Capítulo 6. El período de estancamiento y declinación (1973-1989)
       \begin{enumerate}
       \item La vuelta del peronismo
       \item El gobierno militar
         \item El radicalismo
      \end{enumerate}
     \item Temas
       \begin{enumerate}
       \item El trienio peronista: el plan anti-inflacionario 1973-1975
       \item La política económica del gobierno militar y la situación
         de la economía en 1976.
       \item Crisis de la deuda y decada de los 80. Programas anti-inflacionarios
         \end{enumerate}
         \end{enumerate}
  \end{frame}


 \begin{frame}\frametitle{Programa abreviado: Capítulo 7}
     \begin{enumerate}
     \item Capítulo 7. La decada de los 90 
       \begin{enumerate}
       \item El primer gobierno de Menem
             \end{enumerate}
     \item Temas
       \begin{enumerate}
       \item Las transformaciones económicas de los 90.
         \item La Convertibilidad en persepctiva comparada con la Caja
           de Conversión
         \end{enumerate}
         \end{enumerate}
  \end{frame}
  

\end{document}
  